\documentclass[12pt, a4paper, twoside]{book}

% ─── Packages ───────────────────────────────────────────────────────────────
\usepackage[a4paper, inner=3.0cm, outer=2.5cm, top=2.8cm, bottom=2.8cm, includeheadfoot, headheight=24pt]{geometry}
\usepackage[T1]{fontenc}
\usepackage[utf8]{inputenc}
\usepackage{amsmath, amsthm}
\usepackage[charter, cal=cmcal]{mathdesign}
\usepackage[scaled=0.92]{helvet}
\usepackage{mathtools}
\usepackage{physics}
\usepackage{graphicx}
\usepackage{xcolor}
\usepackage{tikz}
\usetikzlibrary{arrows.meta, decorations.pathmorphing, shapes.geometric,
                positioning, calc, patterns, backgrounds, decorations.markings}
\usepackage{pgfplots}
\pgfplotsset{compat=1.18}
\usepackage{tcolorbox}
\tcbuselibrary{skins, breakable, theorems}
\usepackage{enumitem}
\usepackage{booktabs}
\usepackage{array}
\usepackage{multirow}
\usepackage{fancyhdr}
\usepackage{titlesec}
\usepackage{hyperref}
\usepackage{wrapfig}
\usepackage{float}
\usepackage{caption}
\usepackage{subcaption}
\usepackage{epigraph}
\usepackage{lettrine}
\usepackage{mdframed}
\usepackage{setspace}
\usepackage{parskip}
\usepackage{siunitx}

% ─── Colours ────────────────────────────────────────────────────────────────
\definecolor{deep}{HTML}{0F172A}      % Deep Slate/Charcoal for body text/titles
\definecolor{accent}{HTML}{E11D48}    % Crimson/Coral for problem counters
\definecolor{teal}{HTML}{0D9488}      % Electric Teal for rules and highlights
\definecolor{gold}{HTML}{D97706}      % Gold for puzzle borders
\definecolor{midblue}{HTML}{334155}   % Muted Blue-Gray for section headings
\definecolor{purple}{HTML}{6D28D9}    % Purple for reflection borders
\definecolor{lightbg}{HTML}{F8FAFC}   % Very light gray-blue for exploration bg
\definecolor{warmbg}{HTML}{FEF3C7}    % Warm gold-yellow for puzzle bg
\definecolor{greenbg}{HTML}{F0FDF4}   % Light green for thought/keyidea bg

% ─── Page geometry fix ──────────────────────────────────────────────────────
\setlength{\headheight}{24pt}

% ─── Hyperref ───────────────────────────────────────────────────────────────
\hypersetup{
  colorlinks=true, linkcolor=midblue, urlcolor=accent, citecolor=teal,
  pdftitle={Discovery Physics: Journey 2 — Motion and Its Causes},
  pdfauthor={Discovery Institute},
}

% ─── Page style ─────────────────────────────────────────────────────────────
\pagestyle{fancy}
\fancyhf{}
\fancyhead[LE]{\small\textcolor{deep}{\textit{Discovery Physics}}}
\fancyhead[RO]{\small\textcolor{deep}{\textit{\leftmark}}}
\fancyfoot[C]{\small\textcolor{deep}{\thepage}}
\renewcommand{\headrulewidth}{0.4pt}
\renewcommand{\headrule}{\hbox to\headwidth{\color{teal}\leaders\hrule height \headrulewidth\hfill}}

% ─── Chapter / Section titles ───────────────────────────────────────────────
\titleformat{\chapter}[display]
  {\normalfont\huge\bfseries\color{deep}}
  {\textcolor{accent}{\chaptertitlename\ \thechapter}}{12pt}{\Huge}
  [\vspace{4pt}\textcolor{teal}{\titlerule[1.5pt]}]

\titleformat{\section}
  {\normalfont\Large\bfseries\color{midblue}}
  {\textcolor{accent}{\thesection}}{1em}{}

\titleformat{\subsection}
  {\normalfont\large\bfseries\color{deep}}
  {\textcolor{teal}{\thesubsection}}{1em}{}

% ─── Custom boxes ───────────────────────────────────────────────────────────
\newtcolorbox{puzzlebox}[1][]{
  enhanced, breakable,
  colback=warmbg, colframe=gold,
  fonttitle=\bfseries\large\color{deep}, title={The Puzzle},
  attach boxed title to top left={yshift=-2mm, xshift=4mm},
  boxed title style={colback=gold, colframe=gold}, #1}

\newtcolorbox{explobox}[1][]{
  enhanced, breakable,
  colback=lightbg, colframe=midblue,
  fonttitle=\bfseries\color{white}, coltitle=white, title={Exploration},
  attach boxed title to top left={yshift=-2mm, xshift=4mm},
  boxed title style={colback=midblue, colframe=midblue}, #1}

\newtcolorbox{discoverybox}[1][]{
  enhanced, breakable,
  colback=white, colframe=accent,
  fonttitle=\bfseries\large\color{white}, coltitle=white, title={Discovery},
  attach boxed title to top left={yshift=-2mm, xshift=4mm},
  boxed title style={colback=accent, colframe=accent}, drop shadow, #1}

\newtcolorbox{namebox}[1][]{
  enhanced,
  colback=deep!5, colframe=deep,
  fonttitle=\bfseries\color{white}, coltitle=white, title={What This Is Called},
  attach boxed title to top left={yshift=-2mm, xshift=4mm},
  boxed title style={colback=deep, colframe=deep}, #1}

\newtcolorbox{reflectionbox}[1][]{
  enhanced, breakable,
  colback=purple!5, colframe=purple,
  fonttitle=\bfseries\color{white}, coltitle=white, title={Reflection}, #1}

\newtcolorbox{historybox}[1][]{
  enhanced, breakable,
  colback=gold!8, colframe=gold!70!brown,
  fonttitle=\bfseries\color{deep}, title={Historical Note}, #1}

% ─── Problem counter ────────────────────────────────────────────────────────
\newcounter{problem}[chapter]
\newcommand{\problem}[1][]{%
  \refstepcounter{problem}%
  \medskip\noindent
  \textcolor{accent}{\textbf{Problem \theproblem.}}\ %
  \ifx&#1&\else\textit{(#1)}\ \fi}

\newcommand{\keyidea}[1]{%
  \begin{tcolorbox}[enhanced, colback=teal!8, colframe=teal,
    left=4mm, fontupper=\bfseries\color{deep}]#1\end{tcolorbox}}

% ─── Theorem environments ───────────────────────────────────────────────────
\newtheoremstyle{disc}{8pt}{8pt}{\itshape}{}{\bfseries\color{midblue}}{.}{.5em}{}
\theoremstyle{disc}
\newtheorem{principle}{Principle}[chapter]
\newtheorem{question}{Question}[chapter]

% ─── Spacing ────────────────────────────────────────────────────────────────
\setlength{\epigraphwidth}{0.65\textwidth}
\renewcommand{\epigraphflush}{center}
\setlength{\parskip}{6pt}
\setlength{\parindent}{0pt}
\onehalfspacing

% ════════════════════════════════════════════════════════════════════════════
\begin{document}


\setcounter{chapter}{1}
\setcounter{section}{2}

\section{Module 2.3 — Two Directions at Once}
\label{sec:projectiles}

\epigraph{
  \itshape The laws of nature are but the mathematical thoughts of God.
}{--- Euclid (attributed)}

\subsection{The Central Question}

\begin{tcolorbox}[colback=deep!5, colframe=deep, fontupper=\large\bfseries\color{deep}]
When a ball is thrown sideways, it falls at the same rate as one dropped straight down. Why? And what does this tell us about how forces act?
\end{tcolorbox}

\subsection{The Independence of Motions}

\begin{figure}[H]
  \centering
  \includegraphics[width=0.9\textwidth]{figures/fig_j2_projectile}
  \caption{Left: projectile paths at various launch angles (same initial speed). The maximum range occurs at $45°$. Right: the decomposition of projectile motion into independent horizontal (constant) and vertical (accelerating) components.}
  \label{fig:projectile}
\end{figure}

Drop a ball from a table. At exactly the same moment, throw another ball horizontally off the table at high speed. Which lands first?

The answer, which is counter-intuitive to most people, is that they land simultaneously. The horizontal motion---fast or slow---has absolutely no effect on the vertical fall. The two components of motion are completely independent.

This is Galileo's great insight about projectile motion: a thrown object simultaneously undergoes two independent motions:
\begin{enumerate}[leftmargin=2em]
  \item \textbf{Horizontal}: constant velocity (no force acts horizontally, in the absence of air resistance)
  \item \textbf{Vertical}: uniformly accelerating downward at $g$
\end{enumerate}

\begin{explobox}
\problem A ball is launched horizontally from a height of $\SI{20}{m}$ with speed $\SI{15}{m/s}$.
\begin{enumerate}[label=(\alph*)]
  \item How long does it take to reach the ground?
  \item How far horizontally does it travel?
  \item What is its speed just before hitting the ground? (Both components contribute.)
\end{enumerate}

\problem A projectile is launched at angle $\theta$ above horizontal with initial speed $v_0$. The initial velocity components are $v_{0x} = v_0\cos\theta$ and $v_{0y} = v_0\sin\theta$.
\begin{enumerate}[label=(\alph*)]
  \item Find the time of flight (when does $y = 0$ again?).
  \item Find the horizontal range.
  \item At what angle is the range maximised? Why does this make sense?
\end{enumerate}

\problem A rifle is fired horizontally at a target $\SI{200}{m}$ away. The bullet leaves the barrel at $\SI{600}{m/s}$. By how much does the bullet drop by the time it reaches the target? Is this significant for marksmanship?

\problem \textit{The monkey and the hunter.} A monkey hangs from a branch directly above a point on the ground. A hunter aims the rifle directly at the monkey and fires. At the instant the bullet leaves the rifle, the monkey drops from the branch. Does the bullet hit the monkey, pass above it, or pass below it? Justify your answer using the independence of motions.
\end{explobox}

\subsection{Circular Motion: A Special Case}

An object moving in a circle at constant speed has a changing velocity---because direction is changing. Changing velocity means there is an acceleration, even at constant speed. This acceleration points toward the centre of the circle.

For circular motion with speed $v$ and radius $r$:
\begin{equation}
  a_{\text{centripetal}} = \frac{v^2}{r}
\end{equation}

\begin{explobox}
\problem A car rounds a circular bend of radius $\SI{50}{m}$ at $\SI{60}{km/h}$. What is the centripetal acceleration? Express it as a multiple of $g$.

\problem The Moon orbits Earth at a distance of $\SI{384000}{km}$ with a period of $27.3$ days.
\begin{enumerate}[label=(\alph*)]
  \item What is the Moon's orbital speed?
  \item What is the centripetal acceleration of the Moon?
\end{enumerate}
(Keep your answer---you will need it in Module 2.6.)

\problem A ball on a string is swung in a horizontal circle. If the string is made twice as long but the speed stays the same, how does the centripetal acceleration change? If the speed is doubled but the radius stays the same?
\end{explobox}

% ────────────────────────────────────────────────────────────────────────────

\end{document}
